\noindent\textbf{Elección del procedimiento.} Ahora se mantiene fuerza de tracción $F_t$, no potencia. Los equilibrios son $F_t=S_g+k_DV_1^2$ y $F_t=-S_g+k_DV_2^2$. Restar da $k_D=2S_g/(V_2^2-V_1^2)$; la potencia difiere porque $P_i=F_tV_i$.

\begin{enumerate}

\item Calcular la misma componente del peso.

\[\fitmath{\mathrm{Sg}=\frac{M\cdot g\cdot \mathrm{slope}}{\sqrt{1+{\mathrm{slope}}^{2}}}}\]

\[\fitmath{\mathrm{Sg}=\frac{(1600)\cdot (9.81)\cdot (0.03)}{\sqrt{1+{(0.03)}^{2}}}\approx 470.6682\,\mathrm{N}}\]

\item Eliminar la fuerza común mediante resta de balances.

\[\fitmath{\mathrm{kD}=\frac{2\cdot \mathrm{Sg}}{{V2}^{2}-{V1}^{2}}}\]

\[\fitmath{\mathrm{kD}=\frac{2\cdot (470.6682469)}{{(50.7)}^{2}-{(32.8)}^{2}}\approx 0.629804\,\mathrm{kg/m}}\]

\item Recuperar tracción constante en la subida.

\[\fitmath{\mathrm{Ft}=\mathrm{Sg}+\mathrm{kD}\cdot {V1}^{2}}\]

\[\fitmath{\mathrm{Ft}=(470.6682469)+(0.6298039634)\cdot {(32.8)}^{2}\approx 1148.237\,\mathrm{N}}\]

\item La potencia de subida depende de su velocidad.

\[\fitmath{P1=\mathrm{Ft}\cdot V1}\]

\[\fitmath{P1=(1148.236543)\cdot (32.8)\approx 37662.16\,\mathrm{W}}\]

\item La bajada demanda otra potencia aunque la fuerza sea la misma.

\[\fitmath{P2=\mathrm{Ft}\cdot V2}\]

\[\fitmath{P2=(1148.236543)\cdot (50.7)\approx 58215.59\,\mathrm{W}}\]

\item Despejar el coeficiente de resistencia.

\[\fitmath{C_D=\frac{2\cdot \mathrm{kD}}{\rho\cdot A}}\]

\[\fitmath{C_D=\frac{2\cdot (0.6298039634)}{(1.2)\cdot (2.3)}\approx 0.4563797\,}\]

\end{enumerate}
